
\Data\ID{1003030101}
\Updated{2020-07-15 03:07:12}
\Level{A}
\SubArea{Exponential equations and inequalities}
\LANGen{
\Question{
Solve.
\[ 2^x=-4\]}
{The equation has no solution.}{\( x=2 \)}{\( x=-2 \)}{\( x=-\frac12 \)}{}{}}
\LANGpl{
\Question{
Rozwiąż.
\[ 2^x=-4\]}
{Równanie nie ma rozwiązania.}{\( x=2 \)}{\( x=-2 \)}{\( x=-\frac12 \)}{}{}}
\LANGcs{
\Question{
Řešte danou rovnici.
\[ 2^x=-4\]}
{Rovnice nemá řešení.}{\( x=2 \)}{\( x=-2 \)}{\( x=-\frac12 \)}{}{}}
\LANGsk{
\Question{
Riešte danú rovnicu.
\[ 2^x=-4\]}
{Rovnica nemá riešenie.}{\( x=2 \)}{\( x=-2 \)}{\( x=-\frac12 \)}{}{}}
\LANGes{
\Question{
Resuelve.
\[ 2^x=-4\]}
{La ecuación no tiene solución.}{\( x=2 \)}{\( x=-2 \)}{\( x=-\frac12 \)}{}{}}
\End

\Data\ID{1003030102}
\Updated{2020-07-15 04:07:14}
\Level{A}
\SubArea{Exponential equations and inequalities}
\LANGen{
\Question{
Solve.
\[ 5^{8-2x}=1\]}
{\( x=4 \)}{\( x=-4 \)}{\( x=\frac72 \)}{The equation has no solution.}{}{}}
\LANGpl{
\Question{
Rozwiąż.
\[ 5^{8-2x}=1\]}
{\( x=4 \)}{\( x=-4 \)}{\( x=\frac72 \)}{Równanie nie ma rozwiązania.}{}{}}
\LANGcs{
\Question{
Řešte danou rovnici.
\[ 5^{8-2x}=1\]}
{\( x=4 \)}{\( x=-4 \)}{\( x=\frac72 \)}{Rovnice nemá řešení.}{}{}}
\LANGsk{
\Question{
Riešte danú rovnicu.
\[ 5^{8-2x}=1\]}
{\( x=4 \)}{\( x=-4 \)}{\( x=\frac72 \)}{Rovnica nemá riešenie.}{}{}}
\LANGes{
\Question{
Resuelve.
\[ 5^{8-2x}=1\]}
{\( x=4 \)}{\( x=-4 \)}{\( x=\frac72 \)}{La ecuación no tiene solución.}{}{}}
\End

\Data\ID{1003030103}
\Updated{2020-07-15 03:07:12}
\Level{A}
\SubArea{Exponential equations and inequalities}
\LANGen{
\Question{
Solve.
\[ 8^{2x}=16\sqrt[3]4\]}
{\( x=\frac79 \)}{\( x=\frac{11}{12} \)}{\( x=\frac49 \)}{\( x=\frac13 \)}{}{}}
\LANGpl{
\Question{
Rozwiąż.
\[ 8^{2x}=16\sqrt[3]4\]}
{\( x=\frac79 \)}{\( x=\frac{11}{12} \)}{\( x=\frac49 \)}{\( x=\frac13 \)}{}{}}
\LANGcs{
\Question{
Řešte danou rovnici.
\[ 8^{2x}=16\sqrt[3]4\]}
{\( x=\frac79 \)}{\( x=\frac{11}{12} \)}{\( x=\frac49 \)}{\( x=\frac13 \)}{}{}}
\LANGsk{
\Question{
Riešte danú rovnicu.
\[ 8^{2x}=16\sqrt[3]4\]}
{\( x=\frac79 \)}{\( x=\frac{11}{12} \)}{\( x=\frac49 \)}{\( x=\frac13 \)}{}{}}
\LANGes{
\Question{
Resuelve.
\[ 8^{2x}=16\sqrt[3]4\]}
{\( x=\frac79 \)}{\( x=\frac{11}{12} \)}{\( x=\frac49 \)}{\( x=\frac13 \)}{}{}}
\End

\Data\ID{1003030104}
\Updated{2021-04-12 09:04:31}
\Level{A}
\SubArea{Exponential equations and inequalities}
\LANGen{
\Question{
Decide which of the following intervals contains \( d \), if \( \left(0.1^{d-1}\right)^3=10^{d-1} \).}
{\( [-1;3] \)}{\( [9;13] \)}{\( [4;8] \)}{\( [-6;-2] \)}{}{}}
\LANGpl{
\Question{
Zdecyduj, w którym z podanych przedziałów mieści się \( d \), jeśli \( \left(0{,}1^{d-1}\right)^3=10^{d-1} \).}
{\( [-1;3] \)}{\( [9;13] \)}{\( [4;8] \)}{\( [-6;-2] \)}{}{}}
\LANGcs{
\Question{
Vyberte, který z následujících intervalů obsahuje \( d \), když \( \left(0{,}1^{d-1}\right)^3=10^{d-1} \).}
{\( [-1;3] \)}{\( [9;13] \)}{\( [4;8] \)}{\( [-6;-2] \)}{}{}}
\LANGsk{
\Question{
Vyberte, ktorý z nasledujúcich intervalov obsahuje \( d \), ak \( \left(0{,}1^{d-1}\right)^3=10^{d-1} \).}
{\( [-1;3] \)}{\( [9;13] \)}{\( [4;8] \)}{\( [-6;-2] \)}{}{}}
\LANGes{
\Question{
Decide cuál de los siguientes intervalos contiene a \( d \) si \( \left(0.1^{d-1}\right)^3=10^{d-1} \).}
{\( [-1;3] \)}{\( [9;13] \)}{\( [4;8] \)}{\( [-6;-2] \)}{}{}}
\End

\Data\ID{1003030105}
\Updated{2021-04-13 07:04:51}
\Level{A}
\SubArea{Exponential equations and inequalities}
\LANGen{
\Question{
How many solutions does the following equation have? 
\[ \sqrt{81^{x+1}}=3^{2x}\]}
{No solutions}{Infinitely many solutions}{Exactly one solution}{Exactly two solutions}{}{}}
\LANGpl{
\Question{
Ile rozwiązań ma następujące równanie?
\[ \sqrt{81^{x+1}}=3^{2x}\]}
{Nie ma rozwiązania.}{Ma nieskończenie wiele rozwiązań.}{Ma dokładnie jedno rozwiązanie.}{Ma dokładnie dwa rozwiązania.}{}{}}
\LANGcs{
\Question{
Kolik řešení má následující rovnice? 
\[ \sqrt{81^{x+1}}=3^{2x}\]}
{Žádné řešení}{Nekonečně mnoho řešení}{Právě jedno řešení}{Právě dvě řešení}{}{}}
\LANGsk{
\Question{
Koľko riešení má nasledujúca rovnica? 
\[ \sqrt{81^{x+1}}=3^{2x}\]}
{Žiadne riešenie}{Nekonečne veľa riešení}{Práve jedno riešenie}{Práve dve riešenia}{}{}}
\LANGes{
\Question{
¿Cuántas soluciones tiene la siguiente ecuación?
\[ \sqrt{81^{x+1}}=3^{2x}\]}
{No tiene solución}{Infinitas soluciones}{Una única solución}{Exactamente dos soluciones}{}{}}
\End

\Data\ID{1003030106}
\Updated{2021-04-12 09:04:13}
\Level{A}
\SubArea{Exponential equations and inequalities}
\LANGen{
\Question{
How many solutions does the following equation have?
\[ 2^x\cdot5^x=100^{3-x} \]}
{Exactly one solution}{No solutions}{Infinitely many solutions}{Exactly two solutions}{}{}}
\LANGpl{
\Question{
Ile rozwiązań ma podane równanie?
\[ 2^x\cdot5^x=100^{3-x} \]}
{Dokładnie jedno rozwiązanie}{Brak rozwiązania}{Nieskończenie wiele rozwiązań}{Dokładnie dwa rozwiązania}{}{}}
\LANGcs{
\Question{
Kolik řešení má následující rovnice?
\[ 2^x\cdot5^x=100^{3-x} \]}
{Právě jedno řešení}{Nemá řešení}{Nekonečně mnoho řešení}{Právě dvě řešení}{}{}}
\LANGsk{
\Question{
Koľko riešení má nasledujúca rovnica?
\[ 2^x\cdot5^x=100^{3-x} \]}
{Práve jedno riešenie}{Nemá riešenie}{Nekonečne veľa riešení}{Práve dve riešenia}{}{}}
\LANGes{
\Question{
¿Cuántas soluciones tiene la siguiente ecuación? 
\[ 2^x\cdot5^x=100^{3-x} \]}
{una única solución.}{No tiene solución.}{Un número infinito de soluciones.}{Exactamente dos soluciones.}{}{}}
\End

\Data\ID{1003030107}
\Updated{2020-07-15 03:07:12}
\Level{A}
\SubArea{Exponential equations and inequalities}
\LANGen{
\Question{
Solve.
\[ \sqrt[2x+6]{9^{x+2}}=\sqrt[6]{243}\]}
{\( x=3 \)}{\( x=-8 \)}{\( x=\frac{-24}7 \)}{\( x=-3 \)}{}{}}
\LANGpl{
\Question{
Rozwiąż.
\[ \sqrt[2x+6]{9^{x+2}}=\sqrt[6]{243}\]}
{\( x=3 \)}{\( x=-8 \)}{\( x=\frac{-24}7 \)}{\( x=-3 \)}{}{}}
\LANGcs{
\Question{
Řešte danou rovnici.
\[ \sqrt[2x+6]{9^{x+2}}=\sqrt[6]{243}\]}
{\( x=3 \)}{\( x=-8 \)}{\( x=\frac{-24}7 \)}{\( x=-3 \)}{}{}}
\LANGsk{
\Question{
Riešte danú rovnicu.
\[ \sqrt[2x+6]{9^{x+2}}=\sqrt[6]{243}\]}
{\( x=3 \)}{\( x=-8 \)}{\( x=\frac{-24}7 \)}{\( x=-3 \)}{}{}}
\LANGes{
\Question{
Resuelve.
\[ \sqrt[2x+6]{9^{x+2}}=\sqrt[6]{243}\]}
{\( x=3 \)}{\( x=-8 \)}{\( x=\frac{-24}7 \)}{\( x=-3 \)}{}{}}
\End

\Data\ID{1003044501}
\Updated{2020-07-15 03:07:12}
\Level{A}
\SubArea{Exponential equations and inequalities}
\LANGen{
\Question{
Identify the solution of the equation.
\[ 2^{x^2-x-3}=8. \]}
{\( x_1=3 \), \( x_2=-2 \)}{\( x_1=-3\), \( x_2=2 \)}{\( x_1=0 \), \( x_2=1 \)}{\( x_1=0 \), \( x_2=-1 \)}{}{}}
\LANGpl{
\Question{
Wyznacz rozwiązanie równania.
\[ 2^{x^2-x-3}=8\]}
{\( x_1=3 \), \( x_2=-2 \)}{\( x_1=-3\), \( x_2=2 \)}{\( x_1=0 \), \( x_2=1 \)}{\( x_1=0 \), \( x_2=-1 \)}{}{}}
\LANGcs{
\Question{
Určete kořeny dané rovnice.
\[ 2^{x^2-x-3}=8 \]}
{\( x_1=3 \), \( x_2=-2 \)}{\( x_1=-3\), \( x_2=2 \)}{\( x_1=0 \), \( x_2=1 \)}{\( x_1=0 \), \( x_2=-1 \)}{}{}}
\LANGsk{
\Question{
Určte korene danej rovnice.
\[ 2^{x^2-x-3}=8\]}
{\( x_1=3 \), \( x_2=-2 \)}{\( x_1=-3\), \( x_2=2 \)}{\( x_1=0 \), \( x_2=1 \)}{\( x_1=0 \), \( x_2=-1 \)}{}{}}
\LANGes{
\Question{
Identifica la solución de la ecuación.
\[ 2^{x^2-x-3}=8. \]}
{\( x_1=3 \), \( x_2=-2 \)}{\( x_1=-3\), \( x_2=2 \)}{\( x_1=0 \), \( x_2=1 \)}{\( x_1=0 \), \( x_2=-1 \)}{}{}}
\End

\Data\ID{1003044502}
\Updated{2021-04-13 07:04:39}
\Level{A}
\SubArea{Exponential equations and inequalities}
\LANGen{
\Question{
The following exponential equation has solutions \( x_1 \) and \( x_2 \). Find the sum of \( x_1 \) and \( x_2 \).
\[ \frac{4^{x^2 }}{16^{x+1}} =16\cdot4^x \]}
{\( 3 \)}{\( -3 \)}{\( -1 \)}{\( 2 \)}{}{}}
\LANGpl{
\Question{
Podane równanie wykładnicze ma dwa rozwiązania \( x_1 \) i \( x_2 \). Oblicz sumę \( x_1 \) i \( x_2 \).
\[ \frac{4^{x^2 }}{16^{x+1}} =16\cdot4^x \]}
{\( 3 \)}{\( -3 \)}{\( -1 \)}{\( 2 \)}{}{}}
\LANGcs{
\Question{
Daná exponenciální rovnice má dva kořeny \( x_1 \) a \( x_2 \). Určete součet \( x_1 \) a \( x_2 \).
\[ \frac{4^{x^2 }}{16^{x+1}} =16\cdot4^x \]}
{\( 3 \)}{\( -3 \)}{\( -1 \)}{\( 2 \)}{}{}}
\LANGsk{
\Question{
Daná exponenciálna rovnica má dva korene \( x_1 \) a \( x_2 \). Určte súčet \( x_1 \) a \( x_2 \).
\[ \frac{4^{x^2 }}{16^{x+1}} =16\cdot4^x \]}
{\( 3 \)}{\( -3 \)}{\( -1 \)}{\( 2 \)}{}{}}
\LANGes{
\Question{
La siguiente ecuación tiene como raíces \( x_1 \) y \( x_2 \). Halla la suma de \( x_1 \) y \( x_2 \).
\[ \frac{4^{x^2 }}{16^{x+1}} =16\cdot4^x \]}
{\( 3 \)}{\( -3 \)}{\( -1 \)}{\( 2 \)}{}{}}
\End

\Data\ID{1003044503}
\Updated{2020-07-15 03:07:12}
\Level{A}
\SubArea{Exponential equations and inequalities}
\LANGen{
\Question{
Identify the solution of the equation.
\[ 3^{x-3}-3^{x-2}=2 \]}
{The equation has no solution.}{\( x=3 \)}{\( x=-3 \)}{\( x=2 \)}{}{}}
\LANGpl{
\Question{
Wyznacz rozwiązanie równania.
\[ 3^{x-3}-3^{x-2}=2 \]}
{Równanie nie ma rozwiązania.}{\( x=3 \)}{\( x=-3 \)}{\( x=2 \)}{}{}}
\LANGcs{
\Question{
Najděte řešení dané rovnice.
\[ 3^{x-3}-3^{x-2}=2 \]}
{Rovnice nemá řešení.}{\( x=3 \)}{\( x=-3 \)}{\( x=2 \)}{}{}}
\LANGsk{
\Question{
Nájdite riešenie danej rovnice.
\[ 3^{x-3}-3^{x-2}=2 \]}
{Rovnica nemá riešenie.}{\( x=3 \)}{\( x=-3 \)}{\( x=2 \)}{}{}}
\LANGes{
\Question{
Identifica la solución de la ecuación.
\[ 3^{x-3}-3^{x-2}=2 \]}
{La ecuación no tiene solución.}{\( x=3 \)}{\( x=-3 \)}{\( x=2 \)}{}{}}
\End

\Data\ID{1003044504}
\Updated{2021-04-13 07:04:35}
\Level{A}
\SubArea{Exponential equations and inequalities}
\LANGen{
\Question{
How many solutions does the following equation have?
\[ 5^x+5^{x+1}+5^{x+2}+5^{x+3}=156 \]}
{Exactly one solution}{No solutions}{Infinitely many solutions}{Exactly two solutions}{}{}}
\LANGpl{
\Question{
Ile rozwiązań ma podane równanie?
\[ 5^x+5^{x+1}+5^{x+2}+5^{x+3}=156 \]}
{Dokładnie jedno rozwiązanie}{Brak rozwiązania}{Nieskończenie wiele rozwiązań}{Dokładnie dwa rozwiązania}{}{}}
\LANGcs{
\Question{
Kolik řešení má daná rovnice?
\[ 5^x+5^{x+1}+5^{x+2}+5^{x+3}=156 \]}
{Právě jedno řešení}{Nemá řešení}{Nekonečně mnoho řešení}{Právě dvě řešení}{}{}}
\LANGsk{
\Question{
Koľko riešení má daná rovnica?
\[ 5^x+5^{x+1}+5^{x+2}+5^{x+3}=156 \]}
{Práve jedno riešenie}{Nemá riešenie}{Nekonečne veľa riešení}{Práve dve riešenia}{}{}}
\LANGes{
\Question{
¿Cuántas soluciones tiene la siguiente ecuación? 
\[ 5^x+5^{x+1}+5^{x+2}+5^{x+3}=156 \]}
{Una única solución.}{No tiene ninguna solución.}{Un número infinito de soluciones.}{Exactamente dos soluciones.}{}{}}
\End

\Data\ID{1003044505}
\Updated{2020-07-15 03:07:12}
\Level{A}
\SubArea{Exponential equations and inequalities}
\LANGen{
\Question{
Identify the solution of the equation.
\[ 2^{x-1}-7^{x-3}=2^{x-2}+2^{x-3}\]}
{\( x=3 \)}{The equation has no solution.}{\( x=-3 \)}{\( x=0 \)}{}{}}
\LANGpl{
\Question{
Wyznacz rozwiązanie równania.
\[ 2^{x-1}-7^{x-3}=2^{x-2}+2^{x-3}\]}
{\( x=3 \)}{Równanie nie ma rozwiązania.}{\( x=-3 \)}{\( x=0 \)}{}{}}
\LANGcs{
\Question{
Najděte řešení dané rovnice.
\[ 2^{x-1}-7^{x-3}=2^{x-2}+2^{x-3}\]}
{\( x=3 \)}{Rovnice nemá řešení.}{\( x=-3 \)}{\( x=0 \)}{}{}}
\LANGsk{
\Question{
Nájdite riešenie danej rovnice.
\[ 2^{x-1}-7^{x-3}=2^{x-2}+2^{x-3}\]}
{\( x=3 \)}{Rovnica nemá riešenie.}{\( x=-3 \)}{\( x=0 \)}{}{}}
\LANGes{
\Question{
Identifica la solución de la ecuación.
\[ 2^{x-1}-7^{x-3}=2^{x-2}+2^{x-3}\]}
{\( x=3 \)}{La ecuación no tiene solución.}{\( x=-3 \)}{\( x=0 \)}{}{}}
\End

\Data\ID{1003044506}
\Updated{2020-07-15 03:07:12}
\Level{A}
\SubArea{Exponential equations and inequalities}
\LANGen{
\Question{
Which of the given numbers is the product of all the solutions of the following exponential equations? 
\begin{align*} 3^x+3^{x+2}&=30\\ 3^{x+1}+3^{x-1}&=\frac{10}9 \end{align*}}
{\( -1 \)}{\( 0 \)}{\( 1 \)}{\( 2 \)}{}{}}
\LANGpl{
\Question{
Która z podanych liczb jest iloczynem wszystkich rozwiązań następującego równania wykładniczego?
\begin{align*} 3^x+3^{x+2}&=30\\ 3^{x+1}+3^{x-1}&=\frac{10}9 \end{align*}}
{\( -1 \)}{\( 0 \)}{\( 1 \)}{\( 2 \)}{}{}}
\LANGcs{
\Question{
Které z daných čísel je součinem všech řešení daných exponenciálních rovnic? 
\begin{align*} 3^x+3^{x+2}&=30\\ 3^{x+1}+3^{x-1}&=\frac{10}9 \end{align*}}
{\( -1 \)}{\( 0 \)}{\( 1 \)}{\( 2 \)}{}{}}
\LANGsk{
\Question{
Ktoré z daných čísel je súčinom všetkých riešení daných exponenciálnych rovníc? 
\begin{align*} 3^x+3^{x+2}&=30\\ 3^{x+1}+3^{x-1}&=\frac{10}9 \end{align*}}
{\( -1 \)}{\( 0 \)}{\( 1 \)}{\( 2 \)}{}{}}
\LANGes{
\Question{
¿Cuál de los siguientes números es el producto de todas las soluciones de las siguientes ecuaciones exponenciales?  
\begin{align*} 3^x+3^{x+2}&=30\\ 3^{x+1}+3^{x-1}&=\frac{10}9 \end{align*}}
{\( -1 \)}{\( 0 \)}{\( 1 \)}{\( 2 \)}{}{}}
\End

\Data\ID{1003044507}
\Updated{2021-04-13 07:04:23}
\Level{A}
\SubArea{Exponential equations and inequalities}
\LANGen{
\Question{
How many solutions does the following equation have?
\[ 2^{x+2}+3\cdot3^{x+2}=3^{x+3}+2\cdot2^{x+1} \]}
{Infinitely many solutions}{Exactly one solution}{No solutions}{Exactly two solutions}{}{}}
\LANGpl{
\Question{
Ile rozwiązań ma podane równanie?
\[ 2^{x+2}+3\cdot3^{x+2}=3^{x+3}+2\cdot2^{x+1} \]}
{Nieskończenie wiele rozwiązań}{Dokładnie jedno rozwiązanie}{Brak rozwiązań}{Dokładnie dwa rozwiązania}{}{}}
\LANGcs{
\Question{
Kolik řešení má daná rovnice?
\[ 2^{x+2}+3\cdot3^{x+2}=3^{x+3}+2\cdot2^{x+1} \]}
{Nekonečně mnoho řešení}{Právě jedno řešení}{Nemá řešení}{Právě dvě řešení}{}{}}
\LANGsk{
\Question{
Koľko riešení má daná rovnica?
\[ 2^{x+2}+3\cdot3^{x+2}=3^{x+3}+2\cdot2^{x+1} \]}
{Nekonečne veľa riešení}{Práve jedno riešenie}{Nemá riešenie}{Práve dve riešenia}{}{}}
\LANGes{
\Question{
¿Cuántas soluciones tiene la siguiente ecuación? 
\[ 2^{x+2}+3\cdot3^{x+2}=3^{x+3}+2\cdot2^{x+1} \]}
{Un número infinito de soluciones.}{Una única solución.}{No tiene solución.}{Exactamente dos soluciones.}{}{}}
\End

\Data\ID{1003044601}
\Updated{2020-07-15 03:07:12}
\Level{B}
\SubArea{Exponential equations and inequalities}
\LANGen{
\Question{
Identify the solution of the equation.
\[ 4^{2x}-20\cdot4^x=-64\]}
{\( x_1=1;\ x_2=2 \)}{\( x_1=16;\ x_2=4 \)}{\( x_1=-16;\ x_2=-4 \)}{\( x_1=1;\ x_2=4 \)}{}{}}
\LANGpl{
\Question{
Wyznacz rozwiązanie równania.
\[ 4^{2x}-20\cdot4^x=-64\]}
{\( x_1=1;\ x_2=2 \)}{\( x_1=16;\ x_2=4 \)}{\( x_1=-16;\ x_2=-4 \)}{\( x_1=1;\ x_2=4 \)}{}{}}
\LANGcs{
\Question{
Určete kořeny rovnice.
\[ 4^{2x}-20\cdot4^x=-64\]}
{\( x_1=1;\ x_2=2 \)}{\( x_1=16;\ x_2=4 \)}{\( x_1=-16;\ x_2=-4 \)}{\( x_1=1;\ x_2=4 \)}{}{}}
\LANGsk{
\Question{
Určte korene danej rovnice.
\[ 4^{2x}-20\cdot4^x=-64\]}
{\( x_1=1;\ x_2=2 \)}{\( x_1=16;\ x_2=4 \)}{\( x_1=-16;\ x_2=-4 \)}{\( x_1=1;\ x_2=4 \)}{}{}}
\LANGes{
\Question{
Identifica la solución de la ecuación. 
\[ 4^{2x}-20\cdot4^x=-64\]}
{\( x_1=1;\ x_2=2 \)}{\( x_1=16;\ x_2=4 \)}{\( x_1=-16;\ x_2=-4 \)}{\( x_1=1;\ x_2=4 \)}{}{}}
\End

\Data\ID{1003044602}
\Updated{2021-04-13 07:04:55}
\Level{B}
\SubArea{Exponential equations and inequalities}
\LANGen{
\Question{
Identify the solution of the equation.
\[ 25^x+5=6\cdot5^x \]}
{\( x_1=0;\ x_2=1 \)}{\( x=1 \)}{\( x_1=1;\ x_2=5 \)}{\( x_1=-1;\ x_2=-5 \)}{}{}}
\LANGpl{
\Question{
Wyznacz rozwiązanie równania.
\[ 25^x+5=6\cdot5^x \]}
{\( x_1=0;\ x_2=1 \)}{\( x=1 \)}{\( x_1=1;\ x_2=5 \)}{\( x_1=-1;\ x_2=-5 \)}{}{}}
\LANGcs{
\Question{
Určete kořeny dané rovnice.
\[ 25^x+5=6\cdot5^x \]}
{\( x_1=0;\ x_2=1 \)}{\( x=1 \)}{\( x_1=1;\ x_2=5 \)}{\( x_1=-1;\ x_2=-5 \)}{}{}}
\LANGsk{
\Question{
Určte korene danej rovnice.
\[ 25^x+5=6\cdot5^x \]}
{\( x_1=0;\ x_2=1 \)}{\( x=1 \)}{\( x_1=1;\ x_2=5 \)}{\( x_1=-1;\ x_2=-5 \)}{}{}}
\LANGes{
\Question{
Identifica la solución (o soluciones) de la ecuación.
\[ 25^x+5=6\cdot5^x \]}
{\( x_1=0;\ x_2=1 \)}{\( x=1 \)}{\( x_1=1;\ x_2=5 \)}{\( x_1=-1;\ x_2=-5 \)}{}{}}
\End

\Data\ID{1003044603}
\Updated{2021-04-13 07:04:35}
\Level{B}
\SubArea{Exponential equations and inequalities}
\LANGen{
\Question{
How many solutions does the following equation have?
\[ 81^x-8\cdot9^x=9 \]}
{Exactly one solution}{Exactly two solutions}{No solutions}{Infinitely many solutions}{}{}}
\LANGpl{
\Question{
Ile rozwiązań ma podane równanie?
\[ 81^x-8\cdot9^x=9 \]}
{Dokładnie jedno rozwiązanie}{Dokładnie dwa rozwiązania}{Brak rozwiązań}{Nieskończenie wiele rozwiązań}{}{}}
\LANGcs{
\Question{
Kolik řešení má daná rovnice?
\[ 81^x-8\cdot9^x=9 \]}
{Právě jedno řešení}{Právě dvě řešení}{Nemá řešení}{Nekonečně mnoho řešení}{}{}}
\LANGsk{
\Question{
Koľko riešení má daná rovnica?
\[ 81^x-8\cdot9^x=9 \]}
{Práve jedno riešenie}{Práve dve riešenia}{Nemá riešenie}{Nekonečne veľa riešení}{}{}}
\LANGes{
\Question{
¿Cuántas soluciones hay para la siguiente ecuación? 
\[ 81^x-8\cdot9^x=9 \]}
{Una única solución.}{Exactamente dos soluciones.}{No hay solución.}{Un número infinito de soluciones.}{}{}}
\End

\Data\ID{1003044604}
\Updated{2020-07-15 03:07:12}
\Level{B}
\SubArea{Exponential equations and inequalities}
\LANGen{
\Question{
The following exponential equation has two solutions \( x_1 \) and \( x_2 \). Find the product of \( x_1 \) and \( x_2 \).
\[ 6^x+6^{2-x}=37 \]}
{\( 0 \)}{\( 2 \)}{\( 36 \)}{\( 37 \)}{}{}}
\LANGpl{
\Question{
Podane równanie wykładnicze ma dwa rozwiązania \( x_1 \) i \( x_2 \). Oblicz iloczyn  \( x_1 \) i \( x_2 \).
\[ 6^x+6^{2-x}=37 \]}
{\( 0 \)}{\( 2 \)}{\( 36 \)}{\( 37 \)}{}{}}
\LANGcs{
\Question{
Daná exponenciální rovnice má dva kořeny \( x_1 \) a \( x_2 \). Určete součin \( x_1 \) a \( x_2 \).
\[ 6^x+6^{2-x}=37 \]}
{\( 0 \)}{\( 2 \)}{\( 36 \)}{\( 37 \)}{}{}}
\LANGsk{
\Question{
Daná exponenciálna rovnica má dva korene \( x_1 \) a \( x_2 \). Určte súčin \( x_1 \) a \( x_2 \).
\[ 6^x+6^{2-x}=37 \]}
{\( 0 \)}{\( 2 \)}{\( 36 \)}{\( 37 \)}{}{}}
\LANGes{
\Question{
La siguiente ecuación exponencial tiene dos soluciones \( x_1 \) y \( x_2 \). Halla el producto de \( x_1 \) y \( x_2 \).
\[ 6^x+6^{2-x}=37 \]}
{\( 0 \)}{\( 2 \)}{\( 36 \)}{\( 37 \)}{}{}}
\End

\Data\ID{1003044605}
\Updated{2021-04-13 07:04:03}
\Level{B}
\SubArea{Exponential equations and inequalities}
\LANGen{
\Question{
Identify the solution of the equation.
\[ 3^{2x-2}-244\cdot3^{x-5}=-3^{-3} \]}
{\( x_1=-3;\ x_2=2 \)}{\( x_1=-2;\ x_2=3 \)}{\( x_1=\frac1{27};\ x_2=9 \)}{\( x_1=-9;\ x_2=-\frac1{27} \)}{}{}}
\LANGpl{
\Question{
Wyznacz rozwiązanie równania.
\[ 3^{2x-2}-244\cdot3^{x-5}=-3^{-3} \]}
{\( x_1=-3;\ x_2=2 \)}{\( x_1=-2;\ x_2=3 \)}{\( x_1=\frac1{27};\ x_2=9 \)}{\( x_1=-9;\ x_2=-\frac1{27} \)}{}{}}
\LANGcs{
\Question{
Určete kořeny dané rovnice.
\[ 3^{2x-2}-244\cdot3^{x-5}=-3^{-3} \]}
{\( x_1=-3;\ x_2=2 \)}{\( x_1=-2;\ x_2=3 \)}{\( x_1=\frac1{27};\ x_2=9 \)}{\( x_1=-9;\ x_2=-\frac1{27} \)}{}{}}
\LANGsk{
\Question{
Určte korene danej rovnice.
\[ 3^{2x-2}-244\cdot3^{x-5}=-3^{-3} \]}
{\( x_1=-3;\ x_2=2 \)}{\( x_1=-2;\ x_2=3 \)}{\( x_1=\frac1{27};\ x_2=9 \)}{\( x_1=-9;\ x_2=-\frac1{27} \)}{}{}}
\LANGes{
\Question{
Identifica la solución (o soluciones) de la ecuación..
\[ 3^{2x-2}-244\cdot3^{x-5}=-3^{-3} \]}
{\( x_1=-3;\ x_2=2 \)}{\( x_1=-2;\ x_2=3 \)}{\( x_1=\frac1{27};\ x_2=9 \)}{\( x_1=-9;\ x_2=-\frac1{27} \)}{}{}}
\End

\Data\ID{1003044606}
\Updated{2021-04-13 07:04:28}
\Level{B}
\SubArea{Exponential equations and inequalities}
\LANGen{
\Question{
How many solutions does the following equation have?
\[ 7^{\sqrt x}+7^{1-\sqrt x}=8 \]}
{Exactly two solutions}{Exactly one solution}{No solutions}{Infinitely many solutions}{}{}}
\LANGpl{
\Question{
Ile rozwiązań ma podane równanie?
\[ 7^{\sqrt x}+7^{1-\sqrt x}=8 \]}
{Dokładnie dwa rozwiązania}{Dokładnie jedno rozwiązanie}{Brak rozwiązań}{Nieskończenie wiele rozwiązań}{}{}}
\LANGcs{
\Question{
Kolik řešení má daná rovnice?
\[ 7^{\sqrt x}+7^{1-\sqrt x}=8 \]}
{Právě dvě řešení}{Právě jedno řešení}{Nemá řešení}{Nekonečně mnoho řešení}{}{}}
\LANGsk{
\Question{
Koľko riešení má daná rovnica?
\[ 7^{\sqrt x}+7^{1-\sqrt x}=8 \]}
{Práve dve riešenia}{Práve jedno riešenie}{Nemá riešenie}{Nekonečne veľa riešení}{}{}}
\LANGes{
\Question{
¿Cuántas soluciones tiene la siguiente ecuación? 
\[ 7^{\sqrt x}+7^{1-\sqrt x}=8 \]}
{Exactamente dos soluciones.}{Una única solución.}{No hay solución}{Un número infinito de soluciones.}{}{}}
\End

\Data\ID{1003044607}
\Updated{2021-04-13 07:04:56}
\Level{B}
\SubArea{Exponential equations and inequalities}
\LANGen{
\Question{
Identify the solution of the equation.
\[ 16\cdot16^{\sqrt{2x+5}}-65\cdot4^{\sqrt{2x+5}}+4=0 \]}
{\( x=-2 \)}{\( x_1=-\frac12;\ x_2=2 \)}{\( x_1=\frac1{16};\ x_2=4 \)}{\( x_1=-4;\ x_2=-\frac1{16} \)}{}{}}
\LANGpl{
\Question{
Wyznacz rozwiązanie równania.
\[ 16\cdot16^{\sqrt{2x+5}}-65\cdot4^{\sqrt{2x+5}}+4=0 \]}
{\( x=-2 \)}{\( x_1=-\frac12;\ x_2=2 \)}{\( x_1=\frac1{16};\ x_2=4 \)}{\( x_1=-4;\ x_2=-\frac1{16} \)}{}{}}
\LANGcs{
\Question{
Určete kořeny dané rovnice.
\[ 16\cdot16^{\sqrt{2x+5}}-65\cdot4^{\sqrt{2x+5}}+4=0 \]}
{\( x=-2 \)}{\( x_1=-\frac12;\ x_2=2 \)}{\( x_1=\frac1{16};\ x_2=4 \)}{\( x_1=-4;\ x_2=-\frac1{16} \)}{}{}}
\LANGsk{
\Question{
Určte korene danej rovnice.
\[ 16\cdot16^{\sqrt{2x+5}}-65\cdot4^{\sqrt{2x+5}}+4=0 \]}
{\( x=-2 \)}{\( x_1=-\frac12;\ x_2=2 \)}{\( x_1=\frac1{16};\ x_2=4 \)}{\( x_1=-4;\ x_2=-\frac1{16} \)}{}{}}
\LANGes{
\Question{
Identifica la solución (o soluciones) de la ecuación. 
\[ 16\cdot16^{\sqrt{2x+5}}-65\cdot4^{\sqrt{2x+5}}+4=0 \]}
{\( x=-2 \)}{\( x_1=-\frac12;\ x_2=2 \)}{\( x_1=\frac1{16};\ x_2=4 \)}{\( x_1=-4;\ x_2=-\frac1{16} \)}{}{}}
\End

\Data\ID{1003044701}
\Updated{2021-04-13 07:04:14}
\Level{B}
\SubArea{Exponential equations and inequalities}
\LANGen{
\Question{
Identify the solution of the following equation.
\[10\cdot25^x+10\cdot4^x=29\cdot10^x \]}
{\( x_1=-1;\ x_2=1 \)}{The equation has no solution.}{\( x_1=\frac25;\ x_2=\frac52 \)}{\( x_1=-\frac52;\ x_2=-\frac25 \)}{}{}}
\LANGpl{
\Question{
Wyznacz rozwiązanie podanego równania.
\[10\cdot25^x+10\cdot4^x=29\cdot10^x \]}
{\( x_1=-1;\ x_2=1 \)}{Rozwiązanie nie ma równania.}{\( x_1=\frac25;\ x_2=\frac52 \)}{\( x_1=-\frac52;\ x_2=-\frac25 \)}{}{}}
\LANGcs{
\Question{
Určete kořeny dané rovnice.
\[10\cdot25^x+10\cdot4^x=29\cdot10^x \]}
{\( x_1=-1;\ x_2=1 \)}{Rovnice nemá řešení.}{\( x_1=\frac25;\ x_2=\frac52 \)}{\( x_1=-\frac52;\ x_2=-\frac25 \)}{}{}}
\LANGsk{
\Question{
Určte korene danej rovnice.
\[10\cdot25^x+10\cdot4^x=29\cdot10^x \]}
{\( x_1=-1;\ x_2=1 \)}{Rovnica nemá riešenie.}{\( x_1=\frac25;\ x_2=\frac52 \)}{\( x_1=-\frac52;\ x_2=-\frac25 \)}{}{}}
\LANGes{
\Question{
Identifica la solución (o soluciones) de la ecuación.
\[10\cdot25^x+10\cdot4^x=29\cdot10^x \]}
{\( x_1=-1;\ x_2=1 \)}{The equation has no solution.}{\( x_1=\frac25;\ x_2=\frac52 \)}{\( x_1=-\frac52;\ x_2=-\frac25 \)}{}{}}
\End

\Data\ID{1003044702}
\Updated{2021-04-13 07:04:34}
\Level{B}
\SubArea{Exponential equations and inequalities}
\LANGen{
\Question{
How many solutions does the following equation have?
\[ 2^x\cdot3^{1-x}-2^{1-x}\cdot3^x=1 \]}
{Exactly one solution,  \( x=0 \)}{Exactly one solution, \( x=-1 \)}{Exactly two solutions, \( x_1=-1\), \( x_2=0 \)}{Exactly two solutions, \( x_1=-\frac23\), \( x_2=1 \)}{}{}}
\LANGpl{
\Question{
Ile rozwiązań ma podane równanie?
\[ 2^x\cdot3^{1-x}-2^{1-x}\cdot3^x=1 \]}
{Dokładnie jedno,  \( x=0 \)}{Dokładnie jedno, \( x=-1 \)}{Dokładnie dwa rozwiązania, \( x_1=-1\), \( x_2=0 \)}{Dokładnie dwa rozwiązania, \( x_1=-\frac23\), \( x_2=1 \)}{}{}}
\LANGcs{
\Question{
Kolik řešení má daná rovnice?
\[ 2^x\cdot3^{1-x}-2^{1-x}\cdot3^x=1 \]}
{Právě jedno řešení,  \( x=0 \)}{Právě jedno řešení, \( x=-1 \)}{Právě dvě řešení, \( x_1=-1\), \( x_2=0 \)}{Právě dvě řešení, \( x_1=-\frac23\), \( x_2=1 \)}{}{}}
\LANGsk{
\Question{
Koľko riešení má daná rovnica?
\[ 2^x\cdot3^{1-x}-2^{1-x}\cdot3^x=1 \]}
{Práve jedno riešenie,  \( x=0 \)}{Práve jedno riešenie, \( x=-1 \)}{Práve dve riešenia, \( x_1=-1\), \( x_2=0 \)}{Práve dve riešenia, \( x_1=-\frac23\), \( x_2=1 \)}{}{}}
\LANGes{
\Question{
¿Cuántas soluciones tiene la siguiente ecuación? 
\[ 2^x\cdot3^{1-x}-2^{1-x}\cdot3^x=1 \]}
{Una única solución,  \( x=0 \)}{Una única solución, \( x=-1 \)}{Exactamente dos soluciones, \( x_1=-1\), \( x_2=0 \)}{Exactamente dos soluciones, \( x_1=-\frac23\), \( x_2=1 \)}{}{}}
\End

\Data\ID{1003044703}
\Updated{2021-04-13 07:04:00}
\Level{B}
\SubArea{Exponential equations and inequalities}
\LANGen{
\Question{
The following exponential equation has two solutions \( x_1 \)  and \( x_2 \). Solution \( x_1 \) is a prime number. Find the sum of \( 3x_1 \) and \( 2x_2 \).
\[ 3^{3-x}\cdot4^{x-1}+3^x\cdot4^{2-x}=21 \]}
{\( 8 \)}{\( 7 \)}{\( 0 \)}{\( 3 \)}{}{}}
\LANGpl{
\Question{
Podane równanie wykładnicze ma dwa rozwiązania \( x_1 \)  i \( x_2 \). Rozwiązanie \( x_1 \) jest liczbą pierwszą.  Oblicz sumę  \( 3x_1 \) i \( 2x_2 \).
\[ 3^{3-x}\cdot4^{x-1}+3^x\cdot4^{2-x}=21 \]}
{\( 8 \)}{\( 7 \)}{\( 0 \)}{\( 3 \)}{}{}}
\LANGcs{
\Question{
Daná exponenciální rovnice má dva kořeny \( x_1 \)  a \( x_2 \). Kořen \( x_1 \) je prvočíslo. Určete součet \( 3x_1 \) a \( 2x_2 \).
\[ 3^{3-x}\cdot4^{x-1}+3^x\cdot4^{2-x}=21 \]}
{\( 8 \)}{\( 7 \)}{\( 0 \)}{\( 3 \)}{}{}}
\LANGsk{
\Question{
Daná exponenciálna rovnica má dva korene \( x_1 \)  a \( x_2 \). Koreň \( x_1 \) je prvočíslo. Určte súčet \( 3x_1 \) a \( 2x_2 \).
\[ 3^{3-x}\cdot4^{x-1}+3^x\cdot4^{2-x}=21 \]}
{\( 8 \)}{\( 7 \)}{\( 0 \)}{\( 3 \)}{}{}}
\LANGes{
\Question{
La siguiente ecuación tiene dos soluciones \( x_1 \)  y \( x_2 \). La solución \( x_1 \) es un número primo. Halla la suma de \( 3x_1 \) y \( 2x_2 \).
\[ 3^{3-x}\cdot4^{x-1}+3^x\cdot4^{2-x}=21 \]}
{\( 8 \)}{\( 7 \)}{\( 0 \)}{\( 3 \)}{}{}}
\End

\Data\ID{9000003604}
\Updated{2020-07-15 03:07:34}
\Level{B}
\SubArea{Exponential equations and inequalities}
\LANGen{
\Question{
Solve the following equation. 
\[ 
 10^{x} - 5^{x-1}\cdot 2^{x-2} = 950
\]}
{\(x = 3\)}{\(x = 1\)}{\(x = 2\)}{\(x = 4\)}{}{}}
\LANGpl{
\Question{
Rozwiąż następujące równanie.
\[ 
 10^{x} - 5^{x-1}\cdot 2^{x-2} = 950
\]}
{\(x = 3\)}{\(x = 1\)}{\(x = 2\)}{\(x = 4\)}{}{}}
\LANGcs{
\Question{
Řešením rovnice \(10^{x} - 5^{x-1}\cdot 2^{x-2} = 950\) 
je:}
{\(x = 3\)}{\(x = 1\)}{\(x = 2\)}{\(x = 4\)}{}{}}
\LANGsk{
\Question{
Riešením rovnice \(10^{x} - 5^{x-1}\cdot 2^{x-2} = 950\) 
je koreň:}
{\(x = 3\)}{\(x = 1\)}{\(x = 2\)}{\(x = 4\)}{}{}}
\LANGes{
\Question{
Resuelve la siguiente ecuación.  
\[ 
 10^{x} - 5^{x-1}\cdot 2^{x-2} = 950
\]}
{\(x = 3\)}{\(x = 1\)}{\(x = 2\)}{\(x = 4\)}{}{}}
\End

\Data\ID{9000003608}
\Updated{2020-11-12 04:11:02}
\Level{B}
\SubArea{Exponential equations and inequalities}
\LANGen{
\Question{
Solve the following equation. 
\[ 
 \frac{2}
{3}\cdot 9^{x+1} - 13\cdot 6^{x} + 24\cdot 4^{x-1} = 0
\]}
{\(1,\ -1\)}{\(\frac{3}
{2},\ \frac{2} 
{3}\)}{\(\frac{1}
{2},\ -\frac{1}
{2}\)}{\(\frac{3}
{2},\ -\frac{3}
{2}\)}{}{}}
\LANGpl{
\Question{
Rozwiąż następujące równanie.
\[ 
 \frac{2}
{3}\cdot 9^{x+1} - 13\cdot 6^{x} + 24\cdot 4^{x-1} = 0
\]}
{\(1,\ -1\)}{\(\frac{3}
{2},\ \frac{2} 
{3}\)}{\(\frac{1}
{2},\ -\frac{1}
{2}\)}{\(\frac{3}
{2},\ -\frac{3}
{2}\)}{}{}}
\LANGcs{
\Question{
Najdi všechna řešení následující rovnice.

\[\frac{2}{3}\cdot 9^{x+1} - 13\cdot 6^{x} + 24\cdot 4^{x-1} = 0\]}
{\(1,\ -1\)}{\(\frac{3}
{2},\ \frac{2} 
{3}\)}{\(\frac{1}
{2},\ -\frac{1}
{2}\)}{\(\frac{3}
{2},\ -\frac{3}
{2}\)}{}{}}
\LANGsk{
\Question{
Riešením rovnice \(\frac{2}
{3}\cdot 9^{x+1} - 13\cdot 6^{x} + 24\cdot 4^{x-1} = 0\) 
sú korene:}
{\(1,\ -1\)}{\(\frac{3}
{2},\ \frac{2} 
{3}\)}{\(\frac{1}
{2},\ -\frac{1}
{2}\)}{\(\frac{3}
{2},\ -\frac{3}
{2}\)}{}{}}
\LANGes{
\Question{
Resuelve la siguiente ecuación. 
\[ 
 \frac{2}
{3}\cdot 9^{x+1} - 13\cdot 6^{x} + 24\cdot 4^{x-1} = 0
\]}
{\(1,\ -1\)}{\(\frac{3}
{2},\ \frac{2} 
{3}\)}{\(\frac{1}
{2},\ -\frac{1}
{2}\)}{\(\frac{3}
{2},\ -\frac{3}
{2}\)}{}{}}
\End

\Data\ID{9000003705}
\Updated{2020-11-12 04:11:42}
\Level{B}
\SubArea{Exponential equations and inequalities}
\LANGen{
\Question{
Solve the following exponential equation. 
\[ 
 3^{2x} - 12\cdot 3^{x} + 27 = 0
\]}
{\(x_{1} = 1;\ x_{2} = 2\)}{\(x_{1} = 3;\ x_{2} = 9\)}{\(x_{1} = -1;\ x_{2} = -2\)}{\(x_{1} = -3;\ x_{2} = -9\)}{}{}}
\LANGpl{
\Question{
Rozwiąż następujące równanie wykładnicze.
\[ 
 3^{2x} - 12\cdot 3^{x} + 27 = 0
\]}
{\(x_{1} = 1;\ x_{2} = 2\)}{\(x_{1} = 3;\ x_{2} = 9\)}{\(x_{1} = -1;\ x_{2} = -2\)}{\(x_{1} = -3;\ x_{2} = -9\)}{}{}}
\LANGcs{
\Question{
Najdi všechna řešení následující rovnice. 
\[3^{2x} - 12\cdot 3^{x} + 27 = 0\]}
{\(x_{1} = 1;\ x_{2} = 2\)}{\(x_{1} = 3;\ x_{2} = 9\)}{\(x_{1} = -1;\ x_{2} = -2\)}{\(x_{1} = -3;\ x_{2} = -9\)}{}{}}
\LANGsk{
\Question{
Riešením danej rovnice \(3^{2x} - 12\cdot 3^{x} + 27 = 0\) sú korene:}
{\(x_{1} = 1;\ x_{2} = 2\)}{\(x_{1} = 3;\ x_{2} = 9\)}{\(x_{1} = -1;\ x_{2} = -2\)}{\(x_{1} = -3;\ x_{2} = -9\)}{}{}}
\LANGes{
\Question{
Resuelve la siguiente ecuación exponencial.  
\[ 
 3^{2x} - 12\cdot 3^{x} + 27 = 0
\]}
{\(x_{1} = 1;\ x_{2} = 2\)}{\(x_{1} = 3;\ x_{2} = 9\)}{\(x_{1} = -1;\ x_{2} = -2\)}{\(x_{1} = -3;\ x_{2} = -9\)}{}{}}
\End

\Data\ID{9000003706}
\Updated{2021-04-13 07:04:06}
\Level{B}
\SubArea{Exponential equations and inequalities}
\LANGen{
\Question{
In the following list identify an equation such that neither
\(x = 2\) nor
\(x = -2\) is the
solution of this equation.}
{\(\sqrt{2^{x}}\cdot \sqrt{3^{x}} = 36\)}{\(0.25^{x} = 16\)}{\(6^{-x} = \frac{1} 
{36}\)}{\(25^{x} = \left (\frac{1}
{5}\right )^{x^{2}
 }\)}{}{}}
\LANGpl{
\Question{
Dla którego z poniższych równań rozwiązaniem nie jest ani 
\(x = 2\) ani
\(x = -2\).}
{\(\sqrt{2^{x}}\cdot \sqrt{3^{x}} = 36\)}{\(0.25^{x} = 16\)}{\(6^{-x} = \frac{1} 
{36}\)}{\(25^{x} = \left (\frac{1}
{5}\right )^{x^{2}
 }\)}{}{}}
\LANGcs{
\Question{
Určete, která z daných exponenciálních rovnic nemá řešení ani
\(x = 2\) ani
\(x = -2\).}
{\(\sqrt{2^{x}}\cdot \sqrt{3^{x}} = 36\)}{\(0{,}25^{x} = 16\)}{\(6^{-x} = \frac{1} 
{36}\)}{\(25^{x} = \left (\frac{1}
{5}\right )^{x^{2}
 }\)}{}{}}
\LANGsk{
\Question{
Vyberte rovnicu, ktorej riešením nie je koreň 
\(x = 2\) ani koreň
\(x = -2\).}
{\(\sqrt{2^{x}}\cdot \sqrt{3^{x}} = 36\)}{\(0{,}25^{x} = 16\)}{\(6^{-x} = \frac{1} 
{36}\)}{\(25^{x} = \left (\frac{1}
{5}\right )^{x^{2}
 }\)}{}{}}
\LANGes{
\Question{
Identifica cuál de las siguientes ecuaciones exponenciales no tiene como solución 
\(x = 2\) ni
\(x = -2\).}
{\(\sqrt{2^{x}}\cdot \sqrt{3^{x}} = 36\)}{\(0.25^{x} = 16\)}{\(6^{-x} = \frac{1} 
{36}\)}{\(25^{x} = \left (\frac{1}
{5}\right )^{x^{2}
 }\)}{}{}}
\End

\Data\ID{9000003707}
\Updated{2020-07-15 03:07:34}
\Level{B}
\SubArea{Exponential equations and inequalities}
\LANGen{
\Question{
Each of the following exponential equations has two solutions. Which of the
equations has one positive and one negative solution?}
{\(16^{x} = 0.25^{x^{2}
 }\)}{\(\left (10^{6-x}\right )^{5-x} = 100\)}{\(2^{x^{2}-4x
 } = 1\)}{\(3^{x^{2}-5x+6
 } = 1\)}{}{}}
\LANGpl{
\Question{
Każde z poniższych równań wykładniczych ma dwa rozwiązania. Które z równań ma jedno dodatnie i jedno ujemne rozwiązanie?}
{\(16^{x} = 0.25^{x^{2}
 }\)}{\(\left (10^{6-x}\right )^{5-x} = 100\)}{\(2^{x^{2}-4x
 } = 1\)}{\(3^{x^{2}-5x+6
 } = 1\)}{}{}}
\LANGcs{
\Question{
Následující exponenciální rovnice mají právě dvě řešení. Určete,
která z nich má právě jedno kladné a jedno záporné řešení.}
{\(16^{x} = 0{,}25^{x^{2}
 }\)}{\(\left (10^{6-x}\right )^{5-x} = 100\)}{\(2^{x^{2}-4x
 } = 1\)}{\(3^{x^{2}-5x+6
 } = 1\)}{}{}}
\LANGsk{
\Question{
Každá z nasledujúcich exponenciálnych rovníc má práve dva korene. Určte, ktorá rovnica má práve jeden kladný a jeden záporný koreň.}
{\(16^{x} = 0{,}25^{x^{2}
 }\)}{\(\left (10^{6-x}\right )^{5-x} = 100\)}{\(2^{x^{2}-4x
 } = 1\)}{\(3^{x^{2}-5x+6
 } = 1\)}{}{}}
\LANGes{
\Question{
Cada una de las siguientes ecuaciones exponenciales tiene dos soluciones. ¿Cuál de las ecuaciones tiene una solución positiva y una negativa?}
{\(16^{x} = 0.25^{x^{2}
 }\)}{\(\left (10^{6-x}\right )^{5-x} = 100\)}{\(2^{x^{2}-4x
 } = 1\)}{\(3^{x^{2}-5x+6
 } = 1\)}{}{}}
\End

\Data\ID{9000003708}
\Updated{2021-04-13 07:04:50}
\Level{B}
\SubArea{Exponential equations and inequalities}
\LANGen{
\Question{
Consider the exponential equation 
\[ 
 4^{x+2} - 5\cdot 4^{x+1} + 4^{x-1} + 240 = 0
\]
with \(x\in \mathbb{R}\).
In the following list identify a true statement on this equation.}
{The equation has a unique solution. This solution is a positive integer.}{The equation has a unique solution. This solution is a negative number.}{The equation does not have a solution.}{The equation has two solutions.}{Zero is a
solution of this equation.}{The equation has a unique solution. This solution is a negative integer.}}
\LANGpl{
\Question{
Rozważ podane równanie wykładnicze
\[ 
 4^{x+2} - 5\cdot 4^{x+1} + 4^{x-1} + 240 = 0
\]
z \(x\in \mathbb{R}\).
Które z poniższych zdań dotyczących tego równania jest prawdziwe?}
{Równanie ma tylko jedno rozwiązanie. Tym rozwiązaniem jest liczba całkowita dodatnia.}{Równanie ma tylko jedno rozwiązanie. Tym rozwiązaniem jest liczba ujemna.}{Równanie nie ma rozwiązania.}{Równanie ma dwa rozwiązania.}{Zero jest rozwiązaniem tego równania.}{Równanie ma tylko jedno rozwiązanie. Tym rozwiązaniem jest liczba całkowita ujemna.}}
\LANGcs{
\Question{
Máme dánu exponenciální rovnici
\(4^{x+2} - 5\cdot 4^{x+1} + 4^{x-1} + 240 = 0\) s
neznámou \(x\in \mathbb{R}\).
Vyberte, které z následujících tvrzení je pravdivé.}
{Rovnice má právě jedno řešení
\(x\in \mathbb{N}\).}{Rovnice má právě jedno řešení záporné.}{Rovnice nemá řešení.}{Rovnice má právě dvě řešení.}{Nula je řešením dané rovnice.}{Rovnice má právě jedno řešení
\(x\in \mathbb{Z}^{-}\).}}
\LANGsk{
\Question{
Je daná exponenciálna rovnica
\(4^{x+2} - 5\cdot 4^{x+1} + 4^{x-1} + 240 = 0\) s
neznámou \(x\in \mathbb{R}\).
Vyberte, ktoré z následujúcich tvrdení o rovnici je pravdivé.}
{Rovnica má práve jedno riešenie
\(x\in \mathbb{N}\).}{Rovnica má práve jedno riešenie záporné.}{Rovnica nemá riešenie.}{Rovnica má práve dve riešenia.}{Riešením rovnice je koreň \(x = 0\).}{Rovnica má práve jedno riešenie
\(x\in \mathbb{Z}^{-}\).}}
\LANGes{
\Question{
Dada la ecuación exponencial  
\[ 
 4^{x+2} - 5\cdot 4^{x+1} + 4^{x-1} + 240 = 0
\]
donde \(x\in \mathbb{R}\).
Elige el enunciado verdadero.}
{La ecuación tiene una única solución \(x\in \mathbb{N}\).}{La ecuación tiene una única solución y es un número negativo.}{La ecuación no tiene solución.}{La ecuación tiene dos soluciones.}{La solución de esta ecuación es cero.}{La ecuación tiene única solución \(x\in \mathbb{Z}^{-}\).}}
\End

\Data\ID{1003064101}
\Updated{2020-07-15 03:07:12}
\Level{C}
\SubArea{Exponential equations and inequalities}
\LANGen{
\Question{
Find the solution of the following inequality.
\[ 2^x\geq1 \]}
{\( x\in[0;\infty) \)}{\( x\in(0;\infty) \)}{\( x\in(-\infty;0) \)}{\( x\in(-\infty;0] \)}{}{}}
\LANGpl{
\Question{
Wyznacz zbiór rozwiązań następującej nierówności. 
\[ 2^x\geq1 \]}
{\( x\in\langle0;\infty) \)}{\( x\in(0;\infty) \)}{\( x\in(-\infty;0) \)}{\( x\in(-\infty;0\rangle \)}{}{}}
\LANGcs{
\Question{
Najděte řešení dané nerovnice.
\[ 2^x\geq1 \]}
{\( x\in\langle0;\infty) \)}{\( x\in(0;\infty) \)}{\( x\in(-\infty;0) \)}{\( x\in(-\infty;0\rangle \)}{}{}}
\LANGsk{
\Question{
Určete riešenie danej nerovnice.
\[ 2^x\geq1 \]}
{\( x\in\langle0;\infty) \)}{\( x\in(0;\infty) \)}{\( x\in(-\infty;0) \)}{\( x\in(-\infty;0\rangle \)}{}{}}
\LANGes{
\Question{
Halla la solución para la siguiente inecuación.
\[ 2^x\geq1 \]}
{\( x\in[0;\infty) \)}{\( x\in(0;\infty) \)}{\( x\in(-\infty;0) \)}{\( x\in(-\infty;0] \)}{}{}}
\End

\Data\ID{1003064102}
\Updated{2021-04-13 07:04:31}
\Level{C}
\SubArea{Exponential equations and inequalities}
\LANGen{
\Question{
How many solutions does the following inequality have?
\[ 5^x < 0 \]}
{No solutions}{Infinitely many solutions}{Exactly one solution}{Exactly two solutions}{}{}}
\LANGpl{
\Question{
Ile rozwiązań ma podana nierówność?
\[ 5^x < 0 \]}
{Nie ma rozwiązań}{Nieskończenie wiele rozwiązań}{Dokładnie jedno rozwiązanie}{Dokładnie dwa rozwiązania}{}{}}
\LANGcs{
\Question{
Kolik řešení má daná nerovnice?
\[ 5^x < 0 \]}
{Žádné řešení}{Nekonečně mnoho řešení}{Právě jedno řešení}{Právě dvě řešení}{}{}}
\LANGsk{
\Question{
Koľko riešení má daná nerovnica?
\[ 5^x < 0 \]}
{Žiadne riešenie}{Nekonečne veľa riešení}{Práve jedno riešenie}{Práve dve riešenia}{}{}}
\LANGes{
\Question{
¿Cuántas soluciones tiene la siguiente inecuación? 
\[ 5^x < 0 \]}
{No hay solución.}{Un número infinito de soluciones.}{Una única solución.}{Exactamente dos soluciones.}{}{}}
\End

\Data\ID{1003064103}
\Updated{2020-07-15 03:07:12}
\Level{C}
\SubArea{Exponential equations and inequalities}
\LANGen{
\Question{
Find the solution set of the following inequality.
\[ \left(\frac12\right)^x\geq\left(\frac18\right) \]}
{\( x\in(-\infty;3] \)}{\( x\in(3;\infty) \)}{\( x\in(-\infty;3) \)}{\( x\in[3;\infty) \)}{}{}}
\LANGpl{
\Question{
Wyznacz zbiór rozwiązań podanej nierówności. 
\[ \left(\frac12\right)^x\geq\left(\frac18\right) \]}
{\( x\in(-\infty;3\rangle \)}{\( x\in(3;\infty) \)}{\( x\in(-\infty;3) \)}{\( x\in\langle3;\infty) \)}{}{}}
\LANGcs{
\Question{
Najděte řešení dané nerovnice.
\[ \left(\frac12\right)^x\geq\left(\frac18\right) \]}
{\( x\in(-\infty;3\rangle \)}{\( x\in(3;\infty) \)}{\( x\in(-\infty;3) \)}{\( x\in\langle3;\infty) \)}{}{}}
\LANGsk{
\Question{
Určete riešenie danej nerovnice.
\[ \left(\frac12\right)^x\geq\left(\frac18\right) \]}
{\( x\in(-\infty;3\rangle \)}{\( x\in(3;\infty) \)}{\( x\in(-\infty;3) \)}{\( x\in\langle3;\infty) \)}{}{}}
\LANGes{
\Question{
Halla la solución de la siguiente inecuación. 
\[ \left(\frac12\right)^x\geq\left(\frac18\right) \]}
{\( x\in(-\infty;3] \)}{\( x\in(3;\infty) \)}{\( x\in(-\infty;3) \)}{\( x\in[3;\infty) \)}{}{}}
\End

\Data\ID{1003064104}
\Updated{2021-04-13 07:04:54}
\Level{C}
\SubArea{Exponential equations and inequalities}
\LANGen{
\Question{
How many solutions in the set of positive integers does the following inequality have?
\[ \left(\frac34\right)^{x+1}\leq\left(\frac43\right)^{3-2x} \]}
{\( 4 \)}{\( 5 \)}{\( 1 \)}{\( 0 \)}{}{}}
\LANGpl{
\Question{
Ile rozwiązań w zbiorze liczb całkowitych dodatnich ma podana nierówność? 
\[ \left(\frac34\right)^{x+1}\leq\left(\frac43\right)^{3-2x} \]}
{\( 4 \)}{\( 5 \)}{\( 1 \)}{\( 0 \)}{}{}}
\LANGcs{
\Question{
Kolik řešení v množině přirozených čísel má následující nerovnice?
\[ \left(\frac34\right)^{x+1}\leq\left(\frac43\right)^{3-2x} \]}
{\( 4 \)}{\( 5 \)}{\( 1 \)}{\( 0 \)}{}{}}
\LANGsk{
\Question{
Koľko riešení na množine prirodzených čísel má nasledujúca nerovnica?
\[ \left(\frac34\right)^{x+1}\leq\left(\frac43\right)^{3-2x} \]}
{\( 4 \)}{\( 5 \)}{\( 1 \)}{\( 0 \)}{}{}}
\LANGes{
\Question{
¿Cuántas soluciones de esta inecuación son números naturales? 
\[ \left(\frac34\right)^{x+1}\leq\left(\frac43\right)^{3-2x} \]}
{\( 4 \)}{\( 5 \)}{\( 1 \)}{\( 0 \)}{}{}}
\End

\Data\ID{1003064106}
\Updated{2020-07-15 03:07:12}
\Level{C}
\SubArea{Exponential equations and inequalities}
\LANGen{
\Question{
Find the solution of the following inequality.
\[ 5^x\geq3^x \]}
{\( x\in[0;\infty) \)}{\( x\in[1;\infty) \)}{\( x\in\{0\} \)}{no solution}{}{}}
\LANGpl{
\Question{
Wyznacz zbiór rozwiązań następującej nierówności. 
\[ 5^x\geq3^x \]}
{\( x\in\langle0;\infty) \)}{\( x\in\langle1;\infty) \)}{\( x\in\{0\} \)}{brak rozwiązania}{}{}}
\LANGcs{
\Question{
Najděte řešení dané nerovnice.
\[ 5^x\geq3^x \]}
{\( x\in\langle0;\infty) \)}{\( x\in\langle1;\infty) \)}{\( x\in\{0\} \)}{\( x\in\emptyset \)}{}{}}
\LANGsk{
\Question{
Určete riešenie danej nerovnice.
\[ 5^x\geq3^x \]}
{\( x\in\langle0;\infty) \)}{\( x\in\langle1;\infty) \)}{\( x\in\{0\} \)}{\( x\in\emptyset \)}{}{}}
\LANGes{
\Question{
Halla la solución de la siguiente inecuación.
\[ 5^x\geq3^x \]}
{\( x\in[0;\infty) \)}{\( x\in[1;\infty) \)}{\( x\in\{0\} \)}{No hay solución.}{}{}}
\End

\Data\ID{1003064201}
\Updated{2021-04-13 07:04:57}
\Level{C}
\SubArea{Exponential equations and inequalities}
\LANGen{
\Question{
How many solutions does the following inequality in the set of positive integers have?
\[ 16^x-3\cdot4^x-4\leq0 \]}
{Exactly one solution}{Infinitely many solutions}{Exactly two solutions}{No solutions}{}{}}
\LANGpl{
\Question{
Ile rozwiązań w zbiorze liczb całkowitych dodatnich posiada podana nierówność?
\[ 16^x-3\cdot4^x-4\leq0 \]}
{Dokładnie jedno rozwiązanie}{Nieskończenie wiele rozwiązań}{Dokładnie dwa rozwiązania}{Nie ma rozwiązania}{}{}}
\LANGcs{
\Question{
Kolik řešení v množině přirozených čísel má následující nerovnice?

\[ 16^x-3\cdot4^x-4\leq0 \]}
{Právě jedno řešení}{Nekonečně mnoho řešení}{Právě dvě řešení}{Žádné řešení}{}{}}
\LANGsk{
\Question{
Koľko riešení na množine prirodzených čísel má nasledujúca nerovnica?

\[ 16^x-3\cdot4^x-4\leq0 \]}
{Práve jedno riešenie}{Nekonečne veľa riešení}{Práve dve riešenia}{Žiadne riešenie}{}{}}
\LANGes{
\Question{
¿Cuántas soluciones de la siguiente inecuación hay en el conjunto de números naturales? 
\[ 16^x-3\cdot4^x-4\leq0 \]}
{Una única solución.}{Un número infinito de soluciones.}{Exactamente dos soluciones.}{No hay soluciones.}{}{}}
\End

\Data\ID{1003064202}
\Updated{2020-07-15 03:07:12}
\Level{C}
\SubArea{Exponential equations and inequalities}
\LANGen{
\Question{
Find the solution of the following inequality.
\[ 2^x\leq 3-2^{1-x} \]}
{\( x\in[0; 1] \)}{\( x\in[1; 2] \)}{\( x\in[-2; -1] \)}{No solution}{}{}}
\LANGpl{
\Question{
Wyznacz zbiór rozwiązań następującej nierówności. 
\[ 2^x\leq 3-2^{1-x} \]}
{\( x\in\langle0; 1\rangle \)}{\( x\in\langle1; 2\rangle \)}{\( x\in\langle-2; -1\rangle \)}{Brak rozwiązania}{}{}}
\LANGcs{
\Question{
Řešte danou nerovnici.
\[ 2^x\leq 3-2^{1-x} \]}
{\( x\in\langle0; 1\rangle \)}{\( x\in\langle1; 2\rangle \)}{\( x\in\langle-2; -1\rangle \)}{Nerovnice nemá řešení.}{}{}}
\LANGsk{
\Question{
Riešte danú nerovnicu.
\[ 2^x\leq 3-2^{1-x} \]}
{\( x\in\langle0; 1\rangle \)}{\( x\in\langle1; 2\rangle \)}{\( x\in\langle-2; -1\rangle \)}{Nerovnica nemá riešenie.}{}{}}
\LANGes{
\Question{
Halla la solución de la siguiente inecuación.
\[ 2^x\leq 3-2^{1-x} \]}
{\( x\in[0; 1] \)}{\( x\in[1; 2] \)}{\( x\in[-2; -1] \)}{No hay solución.}{}{}}
\End

\Data\ID{1003064203}
\Updated{2020-07-15 03:07:12}
\Level{C}
\SubArea{Exponential equations and inequalities}
\LANGen{
\Question{
Find the solution of the following inequality.
\[ \sqrt{2^{2x}-2^{3+x} } < 2^x-2^3 \]}
{No solution}{\( x\in(-\infty;3) \)}{\( x\in(8;\infty) \)}{\( x\in[3;\infty) \)}{}{}}
\LANGpl{
\Question{
Wyznacz zbiór rozwiązań następującej nierówności. 
\[ \sqrt{2^{2x}-2^{3+x} } < 2^x-2^3 \]}
{Ta nierówność nie ma rozwiązania.}{\( x\in(-\infty;3) \)}{\( x\in(8;\infty) \)}{\( x\in\langle3;\infty) \)}{}{}}
\LANGcs{
\Question{
Řešte danou nerovnici.
\[ \sqrt{2^{2x}-2^{3+x} } < 2^x-2^3 \]}
{Nerovnice nemá řešení.}{\( x\in(-\infty;3) \)}{\( x\in(8;\infty) \)}{\( x\in\langle3;\infty) \)}{}{}}
\LANGsk{
\Question{
Riešte danú nerovnicu.
\[ \sqrt{2^{2x}-2^{3+x} } < 2^x-2^3 \]}
{Nerovnica nemá riešenie.}{\( x\in(-\infty;3) \)}{\( x\in(8;\infty) \)}{\( x\in\langle3;\infty) \)}{}{}}
\LANGes{
\Question{
Halla la solución de la siguiente inecuación.
\[ \sqrt{2^{2x}-2^{3+x} } < 2^x-2^3 \]}
{No hay solución.}{\( x\in(-\infty;3) \)}{\( x\in(8;\infty) \)}{\( x\in[3;\infty) \)}{}{}}
\End

\Data\ID{1003082601}
\Updated{2020-07-15 03:07:12}
\Level{C}
\SubArea{Exponential equations and inequalities}
\LANGen{
\Question{
Find the set of all solutions of the following inequality.
\[ 27^x-3^{3x}\leq 3 \]}
{\( (-\infty;\infty) \)}{\( (-\infty;0] \)}{\( (-\infty;1] \)}{The inequality has no solution.}{}{}}
\LANGpl{
\Question{
Wyznacz zbiór rozwiązań następującej nierówności. 
\[ 27^x-3^{3x}\leq 3 \]}
{\( (-\infty;\infty) \)}{\( (-\infty;0\rangle \)}{\( (-\infty;1\rangle \)}{Ta nierówność nie ma rozwiązania.}{}{}}
\LANGcs{
\Question{
Vyberte množinu, která je řešením dané nerovnice.
\[ 27^x-3^{3x}\leq 3 \]}
{\( (-\infty;\infty) \)}{\( (-\infty;0\rangle \)}{\( (-\infty;1\rangle \)}{\(\{   \}\)}{}{}}
\LANGsk{
\Question{
Vyberte množinu, ktorá je riešením danej nerovnice.
\[ 27^x-3^{3x}\leq 3 \]}
{\( (-\infty;\infty) \)}{\( (-\infty;0\rangle \)}{\( (-\infty;1\rangle \)}{\(\{   \}\)}{}{}}
\LANGes{
\Question{
Halla el conjunto de todas las soluciones de la siguiente inecuación.
\[ 27^x-3^{3x}\leq 3 \]}
{\( (-\infty;\infty) \)}{\( (-\infty;0] \)}{\( (-\infty;1] \)}{La inecuación no tiene solución.}{}{}}
\End

\Data\ID{1003082602}
\Updated{2020-11-12 04:11:26}
\Level{C}
\SubArea{Exponential equations and inequalities}
\LANGen{
\Question{
How many solutions does the following inequality in the set of positive integers have?
\[ 9^{x+2}-5\cdot9^{x+1} \leq 2916 \]}
{\( 2 \)}{\( 3 \)}{\( 1 \)}{Infinitely many solutions}{}{}}
\LANGpl{
\Question{
Ile rozwiązań w zbiorze liczb całkowitych dodatnich posiada podana nierówność?
\[ 9^{x+2}-5\cdot9^{x+1} \leq 2916 \]}
{\( 2 \)}{\( 3 \)}{\( 1 \)}{Nieskończenie wiele rozwiązań}{}{}}
\LANGcs{
\Question{
Kolik řešení má následující nerovnice v množině přirozených čísel?
\[ 9^{x+2}-5\cdot9^{x+1} \leq 2916 \]}
{\( 2 \)}{\( 3 \)}{\( 1 \)}{Nekonečně mnoho řešení}{}{}}
\LANGsk{
\Question{
Koľko riešení na množine prirodzených čísel má nasledujúca nerovnica?
\[ 9^{x+2}-5\cdot9^{x+1} \leq 2916 \]}
{\( 2 \)}{\( 3 \)}{\( 1 \)}{Nekonečne veľa riešení}{}{}}
\LANGes{
\Question{
¿Cuántas soluciones de la siguiente inecuación hay en el conjunto de números naturales? 
\[ 9^{x+2}-5\cdot9^{x+1} \leq 2916 \]}
{\( 2 \)}{\( 3 \)}{\( 1 \)}{Número infinito de soluciones.}{}{}}
\End

\Data\ID{1003082603}
\Updated{2021-04-13 07:04:39}
\Level{C}
\SubArea{Exponential equations and inequalities}
\LANGen{
\Question{
How many solutions does the following inequality have?
\[ 2^{x-1}-2^{x-3}-2^{x-4} \geq 2^{-4}-2^{-3}-2^{-2} \]}
{Infinitely many solutions}{No solution}{Exactly one solution}{Exactly two solutions}{}{}}
\LANGpl{
\Question{
Ile rozwiązań posiada podana nierówność?
\[ 2^{x-1}-2^{x-3}-2^{x-4} \geq 2^{-4}-2^{-3}-2^{-2} \]}
{Nieskończenie wiele rozwiązań}{Nie ma rozwiązania}{Dokładnie jedno rozwiązanie}{Dokładnie dwa rozwiązania}{}{}}
\LANGcs{
\Question{
Kolik řešení má daná nerovnice?
\[ 2^{x-1}-2^{x-3}-2^{x-4} \geq 2^{-4}-2^{-3}-2^{-2} \]}
{Nekonečně mnoho řešení}{Žádné řešení}{Právě jedno řešení}{Právě dvě řešení}{}{}}
\LANGsk{
\Question{
Koľko riešení má daná nerovnica?
\[ 2^{x-1}-2^{x-3}-2^{x-4} \geq 2^{-4}-2^{-3}-2^{-2} \]}
{Nekonečne veľa riešení}{Žiadne riešenie}{Práve jedno riešenie}{Práve dve riešenia}{}{}}
\LANGes{
\Question{
¿Cuántas soluciones tiene la siguiente inecuación? 
\[ 2^{x-1}-2^{x-3}-2^{x-4} \geq 2^{-4}-2^{-3}-2^{-2} \]}
{Un número infinito de soluciones.}{No hay solución.}{Una única solución.}{Exactamente dos soluciones.}{}{}}
\End

\Data\ID{1003082604}
\Updated{2020-07-15 03:07:12}
\Level{C}
\SubArea{Exponential equations and inequalities}
\LANGen{
\Question{
Find the solution set of the following inequality.
\[ \mathrm{e}^{x+2}+\mathrm{e}\leq\mathrm{e}^{x+1}+\mathrm{e}^2 \]}
{\( (-\infty;0] \)}{\( (-\infty;-1] \)}{\( (-\infty;\infty) \)}{The inequality has no solution.}{}{}}
\LANGpl{
\Question{
Wyznacz zbiór rozwiązań następującej nierówności. 
\[ \mathrm{e}^{x+2}+\mathrm{e}\leq\mathrm{e}^{x+1}+\mathrm{e}^2 \]}
{\( (-\infty;0\rangle \)}{\( (-\infty;-1\rangle \)}{\( (-\infty;\infty) \)}{Ta nierówność nie ma rozwiązania.}{}{}}
\LANGcs{
\Question{
Vyberte množinu, která je řešením dané nerovnice.
\[ \mathrm{e}^{x+2}+\mathrm{e}\leq\mathrm{e}^{x+1}+\mathrm{e}^2 \]}
{\( (-\infty;0\rangle \)}{\( (-\infty;-1\rangle \)}{\( (-\infty;\infty) \)}{\(\{   \}\)}{}{}}
\LANGsk{
\Question{
Vyberte množinu, ktorá je riešením danej nerovnice.
\[ \mathrm{e}^{x+2}+\mathrm{e}\leq\mathrm{e}^{x+1}+\mathrm{e}^2 \]}
{\( (-\infty;0\rangle \)}{\( (-\infty;-1\rangle \)}{\( (-\infty;\infty) \)}{\(\{   \}\)}{}{}}
\LANGes{
\Question{
Halla el conjunto de soluciones de la siguiente inecuación.
\[ \mathrm{e}^{x+2}+\mathrm{e}\leq\mathrm{e}^{x+1}+\mathrm{e}^2 \]}
{\( (-\infty;0] \)}{\( (-\infty;-1] \)}{\( (-\infty;\infty) \)}{La inecuación no tiene solución.}{}{}}
\End

\Data\ID{1003082605}
\Updated{2020-07-15 03:07:12}
\Level{C}
\SubArea{Exponential equations and inequalities}
\LANGen{
\Question{
Solve the following system of two inequalities.
\begin{align*} \left(\frac12\right)^{x+1}-3\left(\frac12\right)^{x+2}+\frac12&\geq0\\ 4^{x+2}-3\cdot4^{x+1} &< 1 \end{align*}}
{The system of inequalities has no solution.}{\( x\in(-\infty;\infty) \)}{\( x\in(-\infty;-1) \)}{\( x\in(-\infty;-1] \)}{}{}}
\LANGpl{
\Question{
Rozwiąż podaną układ nierówności.
\begin{align*} \left(\frac12\right)^{x+1}-3\left(\frac12\right)^{x+2}+\frac12&\geq0\\ 4^{x+2}-3\cdot4^{x+1} &< 1 \end{align*}}
{Układ nierówności nie ma rozwiązania.}{\( x\in(-\infty;\infty) \)}{\( x\in(-\infty;-1) \)}{\( x\in(-\infty;-1\rangle \)}{}{}}
\LANGcs{
\Question{
Řešte danou soustavu nerovnic.
\begin{align*} \left(\frac12\right)^{x+1}-3\left(\frac12\right)^{x+2}+\frac12&\geq0\\ 4^{x+2}-3\cdot4^{x+1} &< 1 \end{align*}}
{Soustava nerovnic nemá řešení.}{\( x\in(-\infty;\infty) \)}{\( x\in(-\infty;-1) \)}{\( x\in(-\infty;-1\rangle \)}{}{}}
\LANGsk{
\Question{
Riešte danú sústavu nerovníc.
\begin{align*} \left(\frac12\right)^{x+1}-3\left(\frac12\right)^{x+2}+\frac12&\geq0\\ 4^{x+2}-3\cdot4^{x+1} &< 1 \end{align*}}
{Sústava nerovníc nemá riešenie.}{\( x\in(-\infty;\infty) \)}{\( x\in(-\infty;-1) \)}{\( x\in(-\infty;-1\rangle \)}{}{}}
\LANGes{
\Question{
Resuelve el siguiente sistema de inecuaciones. 
\begin{align*} \left(\frac12\right)^{x+1}-3\left(\frac12\right)^{x+2}+\frac12&\geq0\\ 4^{x+2}-3\cdot4^{x+1} &< 1 \end{align*}}
{El sistema de inecuaciones no tiene solución.}{\( x\in(-\infty;\infty) \)}{\( x\in(-\infty;-1) \)}{\( x\in(-\infty;-1] \)}{}{}}
\End

\Data\ID{1003082606}
\Updated{2020-07-15 03:07:12}
\Level{C}
\SubArea{Exponential equations and inequalities}
\LANGen{
\Question{
How many of the following inequalities have the same solution set?
\[ \begin{aligned}
2\left(\frac14\right)^{2x-1}-\left(\frac12\right)^{4x-2}-\frac14&\leq 0  \\
2^{4x+4}-15\cdot4^{2x}&\geq 2^4 \\
9^{2x+1}-2\cdot3^5&\geq3^{4x+1}
\end{aligned} \]}
{\( 3 \)}{\( 2 \)}{\( 1 \)}{\( 0 \)}{}{}}
\LANGpl{
\Question{
Ile z poniższych nierówności ma ten sam zbiór rozwiązań?
\[ \begin{aligned}
2\left(\frac14\right)^{2x-1}-\left(\frac12\right)^{4x-2}-\frac14&\leq 0  \\
2^{4x+4}-15\cdot4^{2x}&\geq 2^4 \\
9^{2x+1}-2\cdot3^5&\geq3^{4x+1}
\end{aligned} \]}
{\( 3 \)}{\( 2 \)}{\( 1 \)}{\( 0 \)}{}{}}
\LANGcs{
\Question{
Kolik z uvedených nerovnic má stejné řešení?
\[ \begin{aligned}
2\left(\frac14\right)^{2x-1}-\left(\frac12\right)^{4x-2}-\frac14&\leq 0  \\
2^{4x+4}-15\cdot4^{2x}&\geq 2^4 \\
9^{2x+1}-2\cdot3^5&\geq3^{4x+1}
\end{aligned} \]}
{\( 3 \)}{\( 2 \)}{\( 1 \)}{\( 0 \)}{}{}}
\LANGsk{
\Question{
Dané sú nasledujúce nerovnice. Koľko z nich má rovnaké riešenie?
\[ \begin{aligned}
2\left(\frac14\right)^{2x-1}-\left(\frac12\right)^{4x-2}-\frac14&\leq 0  \\
2^{4x+4}-15\cdot4^{2x}&\geq 2^4 \\
9^{2x+1}-2\cdot3^5&\geq3^{4x+1}
\end{aligned} \]}
{\( 3 \)}{\( 2 \)}{\( 1 \)}{\( 0 \)}{}{}}
\LANGes{
\Question{
¿Cuántas de las siguientes inecuaciones tienen la misma solución? 
\[ \begin{aligned}
2\left(\frac14\right)^{2x-1}-\left(\frac12\right)^{4x-2}-\frac14&\leq 0  \\
2^{4x+4}-15\cdot4^{2x}&\geq 2^4 \\
9^{2x+1}-2\cdot3^5&\geq3^{4x+1}
\end{aligned} \]}
{\( 3 \)}{\( 2 \)}{\( 1 \)}{\( 0 \)}{}{}}
\End

\Data\ID{1003082607}
\Updated{2020-07-15 03:07:12}
\Level{C}
\SubArea{Exponential equations and inequalities}
\LANGen{
\Question{
Find the set of all solutions of the following inequality.
\[ 4^{x+3}-2\cdot5^{x+2} < 5^{x+2}+4^{x+1} \]}
{\( (-1;\infty) \)}{\( (-\infty;-1) \)}{\( (1;\infty) \)}{The inequality has no solution.}{}{}}
\LANGpl{
\Question{
Wyznacz zbiór wszystkich rozwiązań podanej nierówności. 
\[ 4^{x+3}-2\cdot5^{x+2} < 5^{x+2}+4^{x+1} \]}
{\( (-1;\infty) \)}{\( (-\infty;-1) \)}{\( (1;\infty) \)}{Podana nierówność nie ma rozwiązania.}{}{}}
\LANGcs{
\Question{
Vyberte množinu, která je řešením dané nerovnice.
\[ 4^{x+3}-2\cdot5^{x+2} < 5^{x+2}+4^{x+1} \]}
{\( (-1;\infty) \)}{\( (-\infty;-1) \)}{\( (1;\infty) \)}{\(\{   \}\)}{}{}}
\LANGsk{
\Question{
Vyberte množinu, ktorá je riešením danej nerovnice.
\[ 4^{x+3}-2\cdot5^{x+2} < 5^{x+2}+4^{x+1} \]}
{\( (-1;\infty) \)}{\( (-\infty;-1) \)}{\( (1;\infty) \)}{\(\{   \}\)}{}{}}
\LANGes{
\Question{
Halla el conjunto de soluciones de la siguiente inecuación. 
\[ 4^{x+3}-2\cdot5^{x+2} < 5^{x+2}+4^{x+1} \]}
{\( (-1;\infty) \)}{\( (-\infty;-1) \)}{\( (1;\infty) \)}{La inecuación no tiene solución.}{}{}}
\End

\Data\ID{1103064105}
\Updated{2020-11-12 04:11:12}
\Level{C}
\SubArea{Exponential equations and inequalities}
\LANGen{
\Question{
Decide which of the following diagrams  shows the solution set of the given inequality.
\[ 16\cdot4^{2x-3} < 64^{x+1} \]}
{}{}{}{}{}{}}
\LANGpl{
\Question{
Zdecyduj, który wykres przedstawia zbiór rozwiązań podanej nierówności?
\[ 16\cdot4^{2x-3} < 64^{x+1} \]}
{}{}{}{}{}{}}
\LANGcs{
\Question{
Určete, na kterém z následujících obrázků je znázorněna množina řešení dané nerovnice.
\[ 16\cdot4^{2x-3} < 64^{x+1} \]}
{}{}{}{}{}{}}
\LANGsk{
\Question{
Rozhodnite, ktorý z nasledujúcich obrazkov prezentuje riešenie danej nerovnice.
\[ 16\cdot4^{2x-3} < 64^{x+1} \]}
{}{}{}{}{}{}}
\LANGes{
\Question{
NA}
{}{}{}{}{}{}}

\IMGanswer{1103064105a.pdf}
\IMGanswer{1103064105b.pdf}
\IMGanswer{1103064105c.pdf}
\IMGanswer{1103064105d.pdf}\End

\Data\ID{1103064107}
\Updated{2020-11-12 04:11:44}
\Level{C}
\SubArea{Exponential equations and inequalities}
\LANGen{
\Question{
Decide which of the following diagrams  shows the solution set of the given inequality.
\[ \left(\frac13\right)^{x(x+1)} \geq\left(\frac1{27}\right)^2 \]}
{}{}{}{}{}{}}
\LANGpl{
\Question{
Zdecyduj, który wykres przedstawia zbiór rozwiązań podanej nierówności?
\[ \left(\frac13\right)^{x(x+1)} \geq\left(\frac1{27}\right)^2 \]}
{}{}{}{}{}{}}
\LANGcs{
\Question{
Určete, na kterém z následujících obrázků je znázorněna množina řešení dané nerovnice.
\[ \left(\frac13\right)^{x(x+1)} \geq\left(\frac1{27}\right)^2 \]}
{}{}{}{}{}{}}
\LANGsk{
\Question{
Rozhodnite, ktorý z nasledujúcich obrazkov prezentuje riešenie danej nerovnice.
\[ \left(\frac13\right)^{x(x+1)} \geq\left(\frac1{27}\right)^2 \]}
{}{}{}{}{}{}}
\LANGes{
\Question{
NA}
{}{}{}{}{}{}}

\IMGanswer{1103064107a.pdf}
\IMGanswer{1103064107b.pdf}
\IMGanswer{1103064107c.pdf}
\IMGanswer{1103064107d.pdf}\End

\Data\ID{9000003609}
\Updated{2020-07-15 03:07:34}
\Level{C}
\SubArea{Exponential equations and inequalities}
\LANGen{
\Question{
Solve the following inequality. 
\[ 
 \left (\frac{3}
{4}\right )^{x^{2}-2x
 }\leq \frac{4^{x-6}}
{3^{x-6}}
\]}
{\(x\in (-\infty ;-2] \cup [ 3;\infty )\)}{\(x\in \mathbb{R}\setminus \{ - 2;3\}\)}{\(x\in \mathbb{R}\setminus \{ - 3;2\}\)}{\(x\in [ - 2;3] \)}{}{}}
\LANGpl{
\Question{
Rozwiąż następującą nierówność.
\[ 
 \left (\frac{3}
{4}\right )^{x^{2}-2x
 }\leq \frac{4^{x-6}}
{3^{x-6}}
\]}
{\(x\in (-\infty ;-2] \cup [ 3;\infty )\)}{\(x\in \mathbb{R}\setminus \{ - 2;3\}\)}{\(x\in \mathbb{R}\setminus \{ - 3;2\}\)}{\(x\in [ - 2;3] \)}{}{}}
\LANGcs{
\Question{
Řešením nerovnice \(\left (\frac{3}
{4}\right )^{x^{2}-2x
 }\leq \frac{4^{x-6}} 
{3^{x-6}} \) 
jsou čísla:}
{\(x\in (-\infty ;-2\rangle \cup \langle 3;\infty )\)}{\(x\in \mathbb{R}\setminus \{ - 2;3\}\)}{\(x\in \mathbb{R}\setminus \{ - 3;2\}\)}{\(x\in \langle - 2;3\rangle \)}{}{}}
\LANGsk{
\Question{
Riešením nerovnice
 \(\left (\frac{3}
{4}\right )^{x^{2}-2x
 }\leq \frac{4^{x-6}} 
{3^{x-6}} \) 
je:}
{\(x\in (-\infty ;-2\rangle \cup \langle 3;\infty )\)}{\(x\in \mathbb{R}\setminus \{ - 2;3\}\)}{\(x\in \mathbb{R}\setminus \{ - 3;2\}\)}{\(x\in \langle - 2;3\rangle \)}{}{}}
\LANGes{
\Question{
Resuelve la siguiente inecuación. 
\[ 
 \left (\frac{3}
{4}\right )^{x^{2}-2x
 }\leq \frac{4^{x-6}}
{3^{x-6}}
\]}
{\(x\in (-\infty ;-2] \cup [ 3;\infty )\)}{\(x\in \mathbb{R}\setminus \{ - 2;3\}\)}{\(x\in \mathbb{R}\setminus \{ - 3;2\}\)}{\(x\in [ - 2;3] \)}{}{}}
\End

\Data\ID{9000003709}
\Updated{2020-07-15 03:07:34}
\Level{C}
\SubArea{Exponential equations and inequalities}
\LANGen{
\Question{
Solve the following inequality. 
\[ 
 \left (\frac{2}
{3}\right )^{2-3x} < \frac{2^{x+1}} 
{3^{x+1}}
\]}
{\(\left (-\infty ; \frac{1} 
{4}\right )\)}{\(\left (-\frac{1}
{4};\infty \right )\)}{\((-\infty ;4)\)}{\(\left (\frac{1}
{4};\infty \right )\)}{\((4;\infty )\)}{\(\left (-\infty ;-\frac{1}
{4}\right )\)}}
\LANGpl{
\Question{
Rozwiąż podaną nierówność.
\[ 
 \left (\frac{2}
{3}\right )^{2-3x} < \frac{2^{x+1}} 
{3^{x+1}}
\]}
{\(\left (-\infty ; \frac{1} 
{4}\right )\)}{\(\left (-\frac{1}
{4};\infty \right )\)}{\((-\infty ;4)\)}{\(\left (\frac{1}
{4};\infty \right )\)}{\((4;\infty )\)}{\(\left (-\infty ;-\frac{1}
{4}\right )\)}}
\LANGcs{
\Question{
Řešením exponenciální nerovnice
\(\left (\frac{2}
{3}\right )^{2-3x} < \frac{2^{x+1}} 
{3^{x+1}} \) s
neznámou \(x\in \mathbb{R}\) 
je interval:}
{\(\left (-\infty ; \frac{1} 
{4}\right )\)}{\(\left (-\frac{1}
{4};\infty \right )\)}{\((-\infty ;4)\)}{\(\left (\frac{1}
{4};\infty \right )\)}{\((4;\infty )\)}{\(\left (-\infty ;-\frac{1}
{4}\right )\)}}
\LANGsk{
\Question{
Množina všetkých riešení nerovnice 
\(\left (\frac{2}
{3}\right )^{2-3x} < \frac{2^{x+1}} 
{3^{x+1}} \) je:}
{\(\left (-\infty ; \frac{1} 
{4}\right )\)}{\(\left (-\frac{1}
{4};\infty \right )\)}{\((-\infty ;4)\)}{\(\left (\frac{1}
{4};\infty \right )\)}{\((4;\infty )\)}{\(\left (-\infty ;-\frac{1}
{4}\right )\)}}
\LANGes{
\Question{
Resuelve la siguiente inecuación. 
\[ 
 \left (\frac{2}
{3}\right )^{2-3x} < \frac{2^{x+1}} 
{3^{x+1}}
\]}
{\(\left (-\infty ; \frac{1} 
{4}\right )\)}{\(\left (-\frac{1}
{4};\infty \right )\)}{\((-\infty ;4)\)}{\(\left (\frac{1}
{4};\infty \right )\)}{\((4;\infty )\)}{\(\left (-\infty ;-\frac{1}
{4}\right )\)}}
\End
